%!TEX root = ../main.tex
% APÉNDICE D — Registro de prompts utilizados con herramientas de IA
% Se referencia desde preliminares/declaracion-ia.tex (\ref{apx:prompts}).

\chapter{Prompts utilizados con herramientas de IA generativa}\label{apx:prompts}

\begin{guia}
Registro de los prompts más relevantes del trabajo, organizado por
capítulo y tema. Para cada uno: el comando \instruccionIA (qué se le pidió
a la herramienta) y su resultado \respuestaIA (lo que devolvió, ya
revisado y editado por el autor). Se incluyen solo los prompts cuyo
impacto en el documento fue significativo; el registro completo se
conserva en la carpeta compartida con la cátedra (política de uso de IA,
preliminares/declaracion-ia).
\end{guia}

\section{Criterios de registro}

\begin{itemize}
    \item \textbf{Trazabilidad}: cada entrada indica la herramienta y
          versión utilizada y la fecha aproximada.
    \item \textbf{Verificación}: toda cifra, cita o referencia sugerida
          por la herramienta se confirmó en la fuente original antes de
          incorporarla (regla 3 de la declaración de IA).
    \item \textbf{Autoría}: las salidas se transcriben solo como
          evidencia; el texto final del documento puede diferir por la
          revisión y adaptación del autor.
\end{itemize}

\section{Asistencia en el marco teórico y estado del arte}

\instruccionIA[1]{Resumí los argumentos principales de este paper sobre
evaluación de tutores y ranking de compatibilidad en plataformas de
tutoría en línea, en tres párrafos en español.}
\respuestaIA[1]{(resumen devuelto por la herramienta, verificado contra
el texto original)}

\instruccionIA[2]{Compará los modelos de negocio, comisiones y esquemas de
pago de Superprof, Preply, Wyzant y Tus Clases de manera tabular.}
\respuestaIA[2]{(tabla devuelta, contrastada con las tarifas públicas de
cada plataforma antes de usarse en \ref{tab:comparativa})}

\section{Asistencia en la justificación económica}

\instruccionIA[3]{Calculá el VAN a tasa 15\,\% y la TIR del flujo de caja
[Año 0 ... Año 5] de esta tabla de volúmenes y costos.}
\respuestaIA[3]{(cálculos devueltos; se replicaron de forma independiente
en hoja de cálculo antes de incluirlos en \ref{tab:flujo-caja} y en la
sección de indicadores)}

\section{Asistencia en redacción y código}

\instruccionIA[4]{Reescribí este párrafo en registro académico impersonal
sin cambiar su contenido.}
\respuestaIA[4]{(texto devuelto por la herramienta, luego editado)}

\instruccionIA[5]{Generá el scaffolding de una clase Java de Spring Boot
para el endpoint de reserva de una sesión, con validación de permisos por
rol.}
\respuestaIA[5]{(código devuelto, revisado y adaptado por el autor; las
pruebas se escribieron a mano)}